\documentclass[instructions.tex]{subfiles}
\begin{document}
 
\chapter{De la réparation des églises.}
 
\primo Nos églises sont l’honneur et la gloire de la religion ; ce sont des maisons saintes, consacrées à Dieu, des temples augustes où Jésus-Christ habite ; mais elles sont souvent la honte des chrétiens et l’opprobre de la religion.
 
Les appartemens des princes et des riches de la terre sont ornés : tout y brille en or, en magnificence, tandis qu’on voit des églises dans une pauvreté et une indécence qui font horreur. Quelle honte pour tant de seigneurs et de dames, qui font de grandes dépenses en jeux, en équipages, en vaisselle, en parures, en bâtimens, et qui ne donnent pas une obole pour la décoration ou la réparation du lieu saint ; qui peut-être s’y opposent, et n’entretiennent pas même leurs chapelles !
 
On veut avoir à l’église des places honorables, on y place des armoiries, on y fait des deuils superbes et des mausolées, on les entoure de titres et d’écussons ; voilà à quoi on reconnaît un noble ou un riche : mais y reconnaît-on un chrétien ? Tout cela se fait-il par religion, pour glorifier Dieu, ou bien par vanité, et pour faire honneur à sa famille ?
 
On bâtit des églises magnifiques dans les villes ; on le doit : peut-on employer trop de magnificence pour la maison de Dieu ? Mais cette superbe décoration n’est-elle pas plutôt pour embellir une ville que pour honorer Dieu ? Si c’est par un pur motif de religion, pourquoi a-t-on si peu de zèle pour tant d’églises de la campagne, très-pauvres et indécemment ornées ? Ces églises sont bonnes, dit-on, pour des villageois, mais sont-elles bonnes et conviennent-elles au Roi des rois qui les habite ?
 
Après avoir vu les maisons et les palais des grands du monde, on rougit en voyant certaines églises qui servent de palais à Jésus-Christ ! il seroit bien à souhaiter, à l’horreur de le dire, que ces églises fussent aussi magnifiques que les écuries des chevaux des grands seigneurs.
 
Ne reverra-t-on jamais ces heureux temps où les empereurs et les rois se faisaient honneur de travailler de leurs mains pour la maison de Dieu, et où les fidèles, à l’envi, se dépouillaient de ce qu’ils avaient de plus rare pour l’ornement du sanctuaire ?
 
De toutes les dépenses, on ne plaint que celles qu’on fait pour Dieu. Des paroissiens dépensent dans une année, en débauches, en procès, en superfluités, plus qu’il ne faudrait pour réparer et orner une église ; et ils murmurent aussitôt qu’un pasteur leur parle de réparation. Retranchez vos dépenses inutiles, et vous aurez assez pour entretenir la maison de Dieu.
 
Combien de bénéficiers qui tirent de gros revenus de leurs églises, de leurs abbayes, de leurs chapelles, et qui ne donnent rien à celle qui les nourrit ! On laisse à de pauvres peuples des réparations que la conscience exige de nous.
 
On se plaint des calamités et de la misère des temps ; n’en trouverons-nous point la cause dans le peu de zèle qu’on a pour la maison de Dieu ? Vous avez soin de vos maisons, disoit le Seigneur aux Juifs, et vous abandonnez mon temple ; vous le laissez sans réparations ; c’est pour cela que j’ai envoyé la stérilité, la sécheresse, la désolation dans vos campagnes, que j’ai fait périr vos troupeaux et le travail de vos mains\footnote{ Agg. 1.}.
 
\section{Résolutions}
 
\primo Je contribuerai, autant qu’il sera en mon pouvoir, à la réparation des églises.
 
\secundo J’aurai aussi du zèle pour que les églises, les ornemens et les vases sacrés soient tenus proprement et avec décence.
 
\end{document}
